\exercisegroup

\begin{enumerate}
\def\labelenumi{\arabic{enumi})}
\item
  Let P be a proper pointed convex cone, with vertex 0, in E a vector space over R and p be a semi-norm on E and V the set of points \(x \in E\) for which \(p(x) < 1\) . Let M be a vector subspace of E and f be a linear form on M. There exists a linear form g on E, which extends f and is such that it is \(\geqslant 0\) in P and \(|g(x)| \leqslant p(x)\) for all \(x \in E\) , if and only if, for all \(x \in M \cap (V + P)\) , we have \(f(x) > -1\) . (To see that the condition is sufficient consider a point \(x_0 \in M\) such that, \(f(x_0) = 1\) , the cone Q of vertex 0 generated by \(x_0 + V\) , and apply the cor. of the prop. 1 (II, p.~21) to the space E carrying the relation of preorder for which \(P + Q\) is the cone of elements \(\geqslant 0\) .)
\item
  For a set S let \(F = \mathcal{B}(S)\) be the Banach space of the real-valued bounded functions in S (I, p.~4) and let M be a vector subspace of a normed space E. Show that, for every continuous linear mapping f of M in F, there exists a continuous linear mapping g of E in F, that is an extension of f and such that \(\|g\| = \|f\|\) .
\item
  Let E be a vector space over R and let p be a sublinear function on E (II, p.~20). Let A be a convex set such that \(\inf p(y) > -\infty\) .
\end{enumerate}

\begin{enumerate}
\def\labelenumi{\alph{enumi})}
\tightlist
\item
  Show that the function
\end{enumerate}

\[
q (x) = \inf _ {z \in \mathrm{A}, t \geqslant 0} \left(p (x + t z) - t. \inf _ {y \in \mathrm{A}} p (y)\right)
\]

is a sublinear function on E such that \(-p(-x) \leqslant q(x) \leqslant p(x)\) . b) Show that there exists a linear form h on E such that \(h(x) \leqslant p(x)\) in E and that \(\inf_{y \in A} p(y) = \inf_{y \in A} h(y)\) (take h such that \(h(x) \leqslant q(x)\) ).

\begin{enumerate}
\def\labelenumi{\arabic{enumi})}
\setcounter{enumi}{3}
\tightlist
\item
  Let A be a non-empty set of E, a vector space over R, and p a sub-linear function on E. Let B be the set of \(z \in E\) such that \(\inf_{x \in A} p(x - z) \leqslant 0\) ; we have \(A \subset B\) and \(\inf_{x \in A} p(x) \leqslant p(z)\) for all \(z \in B\) ; from which \(\inf_{x \in A} p(x) = \inf_{z \in B} p(z)\) .
\end{enumerate}


\begin{enumerate}
\def\labelenumi{\alph{enumi})}
\item
  Show that the set of the \(y \in \mathbf{E}\) such that \(\inf_{z \in \mathbf{B}} p(z - y) \leqslant 0\) is the set \(\mathbf{B}\) .
\item
  Deduce from a) that the intersection of B and any affine line D in E is closed in D (show that, whatever the points a, b of E, the function \(t \rightarrow p(a + tb)\) is continuous in R).
\item
  Suppose that for each pair of points x, y of A there exists \(z \in A\) such that \(p(z - \frac{1}{2}(x + y)) \leqslant 0\) . Show that, for each pair of points u, v of B we have \(\frac{1}{2}(u + v) \in \mathbf{B}\) (write
\end{enumerate}

\[
z - \frac {1}{2} (u + v) = \left(z - \frac {1}{2} (x + y)\right) + \frac {1}{2} (x - u) + \frac {1}{2} (y - v)
\]

for \(x, y, z\) in A). Deduce that B is then convex (use \(b\) ).

\begin{enumerate}
\def\labelenumi{\alph{enumi})}
\setcounter{enumi}{3}
\tightlist
\item
  Under the hypotheses of \(c\) ), show that there exists a linear form \(h\) on \(\mathbf{E}\) such that \(h(x) \leqslant p(x)\) and that we have \(\inf_{y \in \mathbf{A}} p(y) = \inf_{y \in \mathbf{A}} h(y)\) (use \(c\) ) and exerc. 3).
\end{enumerate}

\begin{enumerate}
\def\labelenumi{\arabic{enumi})}
\setcounter{enumi}{4}
\item
  Let A be a non-empty subset of a vector space E over R and let p be a sublinear function on E. Suppose that, for every pair of points x, y of A, there exists \(z \in A\) such that \(p(z - (x + y)) \leqslant 0\) and that \(p(x) \geqslant 0\) for all \(x \in A\) . Show that there exists a linear form h on E such that \(h(x) \leqslant p(x)\) in E and that \(h(x) \geqslant 0\) for \(x \in A\) . (Apply exerc. 4, c) to the union of the \(\frac{1}{n}A\) for n (integer) \(\geqslant 1\) .)
\item
  \begin{enumerate}
  \def\labelenumii{\alph{enumii})}
  \tightlist
  \item
    Let H be a hyperplane in E a vector space over R and let p be a sublinear function on E. Let f be a linear form on H, such that \(f(y) \leqslant p(y)\) in H. Let a be a point of \(\complement H\), and let h be the linear form on E which extends f and is such that \(h(a) = \inf_{y \in H} (f(y + p(a - y)))\) . Then \(h(x) \leqslant p(x)\) in E. Show that for every linear form \(g\) on E extending \(f\) and such that \(g(x) \leqslant p(x)\) in E, we also have \(g(a) \leqslant h(a)\) .
  \end{enumerate}
\end{enumerate}

\begin{enumerate}
\def\labelenumi{\alph{enumi})}
\setcounter{enumi}{1}
\tightlist
\item
  Let V be a vector subspace of E and f a linear form on V such that \(f(y) \leqslant p(y)\) in V. Let S be a non-empty set of E. Show that there exists a linear form h on E that extends f, such that \(h(x) \leqslant p(x)\) in E and that there is no other linear form g on E extending f such that \(g(x) \leqslant p(x)\) in E that is distinct from h and such that \(g(x) \geqslant h(x)\) in S. (Consider the set F of pairs \((V', f')\) where \(V'\) is a vector subspace containing V and \(f'\) a linear form on \(V'\) extending f and such that \(f'(z) \leqslant p(z)\) in \(V'\) and further such that there is no other linear form \(f''\) on \(V'\) with the same properties and such that \(f''(z) \geqslant f'(z)\) in \(S \cap V'\) . Order F and use a) and th. 2 of S, III, § 2.4.)
\end{enumerate}

\begin{enumerate}
\def\labelenumi{\arabic{enumi})}
\setcounter{enumi}{6}
\tightlist
\item
  Let \(\mathrm{T}\) be a commutative monoid (A, I, § 2.1) carrying a preorder relation \(x \leqslant y\) such that if \(x \leqslant y\) then \(x + z \leqslant y + z\) for all \(z \in \mathrm{T}\) . A mapping \(f\) of \(\mathrm{T}\) in \(\mathbf{R} \cup \{-\infty\}\) is called additive (resp. subadditive, resp. superadditive) if we have
\end{enumerate}

\[
f (x + y) = f (x) + f (y) \quad (\text { resp. } f (x + y) \leqslant f (x) + f (y), \quad \text { resp. } f (x + y) \geqslant f (x) + f (y))
\]

for any \(x, y\) in \(\mathbf{T}\) .

\begin{enumerate}
\def\labelenumi{\alph{enumi})}
\item
  If \(g\) is subadditive and increasing in \(\mathbf{T}\) , then the function \(h(x) = \inf_{n > 0} g(nx) / n\) is subadditive and increasing; we have \(h \leqslant g\) and \(h(0) = 0\) if \(g(0) \geqslant 0\) .
\item
  Under the same hypotheses suppose that there exist two elements \(x_{1}, x_{2}\) of \(T\) and two real numbers \(\xi_{1}, \xi_{2}\) such that \(\xi_{1} < g(x_{1}), \xi_{2} < g(x_{2})\) and \(g(x_{1} + x_{2}) < \xi_{1} + \xi_{2}\) . Let \(y_{1}, y_{2}, z_{1}, z_{2}\) be four elements of \(T\) , let \(n_{1}, n_{2}\) be two integers \(\geqslant 0\) and let \(\alpha_{1}, \alpha_{2}\) be two real numbers such that
\end{enumerate}

\[
\begin{array}{l} n _ {1} \xi_ {1} + g (z _ {1}) <   \alpha_ {1}, y _ {1} \leqslant n _ {1} x _ {1} + z _ {1} \\ n _ {2} \xi_ {2} + g (z _ {2}) <   \alpha_ {2}, y _ {2} \leqslant n _ {2} x _ {2} + z _ {2}. \end{array}
\]

Show that then \(g(n_2y_1 + n_1y_2) < n_2\alpha_1 + n_1\alpha_2\) .

\begin{enumerate}
\def\labelenumi{\alph{enumi})}
\setcounter{enumi}{2}
\tightlist
\item
  Let \(\omega\) be a superadditive function on T such that \(\omega(0) = 0\) , and let \(\Omega\) be an increasing subadditive function on T such that \(\omega(x) \leqslant \Omega(x)\) in T. Show that there exists an increasing additive function \(f\) on T such that \(\omega(x) \leqslant f(x) \leqslant \Omega(x)\) in T. (Remark that the set of increasing subadditive functions \(g\) on T such that \(\omega(x) \leqslant g(x) \leqslant \Omega(x)\) in T is non-empty and inductive for the relation \(\geqslant\) , and take a minimal element of this set for \(f\) ; show using \(a\) ) that \(f(0) = 0\) . To show that there cannot exist pairs of elements of T, \((x_1, x_2)\) such that \(f(x_1 + x_2) < f(x_1) + f(x_2)\) remark that if \(\xi_j \in \mathbf{R}\) , and \(h_j(x) = \inf(n\xi_j + f(y))\) where \(n\) varies in the set of integers \(\geqslant 0\) and \(y\) in the set of elements of T such that \(x \leqslant nx_j + y\) , then \(h_j\) is increasing and subadditive in T ( \(j = 1, 2\) ), \(h_j(x_j) \leqslant \xi_j\) and \(h_j(x) \leqslant f(x)\) for all \(x \in T\) . Then use the definition of \(f\) and part \(b\) ) to obtain a contradiction.)
\end{enumerate}

¶ * 8) a) Let K be a non-discrete valued division ring of which the absolute value is an ultrametric, non linearly compact(cf.~CA, III, § 2, exerc. 15); then there exists a well ordered set I of numbers > 0 and a family (B(ρ)) \(_{ρ \in I}\) of closed balls in K such that the relation ρ < ρ' implies B(ρ) ⊂ B(ρ'), that B(ρ) has radius ρ and that the intersection of the B(ρ) is empty (CA, VI, § 5, exerc. 5). For every x ∈ K, there exists ρ ∈ I such that x ∉ B(ρ); show that the number φ(x) = \(|x-y|\) for a \(y \in \mathbf{B}(\rho)\) depends neither on \(y \in \mathbf{B}(\rho)\) , nor on \(\rho \in \mathbf{I}\) such that \(x \notin \mathbf{B}(\rho)\) . If \(\rho \in \mathbf{I}\) is such that \(x \in \mathbf{B}(\rho)\) then \(\phi(x) \leqslant \rho\) . This being so, for \((x_1, x_2) \in \mathbf{K}^2\) , put \(\| (x_1, x_2) \| = |x_1|\) if \(x_2 = 0\) , and \(\| (x_1, x_2) \| = |x_2| \phi(x_2^{-1} x_1)\) if \(x_2 \neq 0\) . Show that \(\| (x_1, x_2) \|\) is an ultranorm on \(\mathbf{K}^2\) (I, p.~26, exerc. 12) and show that there does not exist any projection of norm 1 of \(\mathbf{K}^2\) on \(\mathbf{K} \times \{0\}\) .

\begin{enumerate}
\def\labelenumi{\alph{enumi})}
\setcounter{enumi}{1}
\item
  Let K be a complete non-discrete valued division ring of which the absolute value is an ultrametric and which is linearly compact. Let E be a vector space of dimension 2 on K with an ultranorm and let D be a line in E; show that for all points \(x \in \mathbf{E}\) , there exists \(y \in \mathbf{D}\) such that \(d(x, D) = d(x, y) = \| x - y \|\) (note that the intersection of D and of a ball of centre \(x\) is a ball in D).
\item
  Deduce from \(a\) ) and \(b\) ) that for a complete non-discrete valued division ring \(K\) , of which the absolute value is an ultrametric, the following properties are equivalent:
\end{enumerate}

\(\alpha)\) K is linearly compact.

β) For every ultranormed vector space E on K, for every vector subspace F of E and every continuous linear form f on F there exists a continuous linear form g on E that extends f and is such that \(\|g\| = \|f\|\) . (Reduce to the case where E is of dimension 2 and use b).) *

